\section{问题背景与问题重述}

热风烘干是中药材产地初加工中最常见的干燥方式，通常包含预热平衡与恒温干燥两个
阶段：先把冷料升温到接近热风温度，再在温湿度基本恒定的环境中持续排出水分。
烘房的送风温度与送风含湿量是唯一可直接调控的工艺量，而药材内部的温度与含水率
并不能实时测量。工艺参数选得不合适，干燥效率、能耗与成品品质会同时受损；
逐一试烘的代价又很高。因此需要建立能刻画药材内部温度场与含水率场演化的
数学模型，用数值仿真代替试验，回答“内部什么时候热透”“什么时候各处都达标干燥”
这两类问题\upcite{khallaf2022herbs,le2025acori}。

\subsection{题设条件}

待烘药材形状近似圆柱，长 $25$ cm、半径 $2$ cm。烘干开始时药材温度为 $28$ °C，
水分浓度（即干基含水率，指每千克绝干物料所含水的千克数）为 $2.55$ kg/kg，
两者在药材内部均匀分布。附件 1 给出烘房热风的温度（单位 °C）与水分浓度
（即空气含湿量，单位 kg/kg）从 $0$ 到 $14400$ s 每 $60$ s 一个点的实测时程。
附件 2 给出烘干过程中各时间点药材的半径（单位 cm），时间跨度 $0$ 至 $259200$ s、
间隔 $1800$ s。附录 2 给出预热平衡阶段所需的密度、比热容、热传导系数、
对流换热系数、对流传质系数与水分扩散系数；附录 3 与附录 4 分别给出另外两套
随含水率与温度变化的物性经验公式。整个烘干过程一般持续 $2$ 至 $3$ 天，
按烘干要求，药材各处的水分浓度应低于 $0.15$ kg/kg。全部结果保留四位小数。

\subsection{待求问题}

\textbf{问题一}　建立预热平衡阶段药材温度和水分浓度变化规律的数学模型，
相关参数取附录 2。给出 $100$、$300$、$600$、$900$、$1200$、$1500$、$1800$ s
七个时刻，到药材中心距离 $0$、$0.5$、$1$、$1.5$、$2$ cm 五个位置处的温度与
水分浓度，并另行提交 $1800$ s 内每隔 $1$ s、到药材中心距离每隔 $0.1$ cm 的
完整结果。

\textbf{问题二}　建立整个烘干过程药材温度和水分浓度变化规律的数学模型，
经验公式统一采用附录 3。给出 $3$ h 内每隔 $0.5$ h、上述五个位置处的温度与
水分浓度，并另行提交每隔 $1$ s、每隔 $0.1$ cm 的完整结果。

\textbf{问题三}　在药材各处水分浓度均低于 $0.15$ kg/kg 的烘干要求下，
确定烘干所需要的时间（单位 h）。给出每隔 $6$ h、到药材中心距离每隔 $0.5$ cm
的水分浓度，末行为烘干结束时刻，并另行提交每隔 $60$ s、每隔 $0.1$ cm 的
完整结果。

\textbf{问题四}　考虑药材因水分流失发生的尺寸变化，根据附件 2 确定烘干时长，
经验公式见附录 4。给出每隔 $6$ h、每隔 $0.5$ cm 的水分浓度，末列为药材表面处的
数值，末行为烘干结束时刻，并另行提交每隔 $60$ s、每隔 $0.1$ cm 的完整结果。
